problem:  
Let \((M^{2n+1},T^{1,0}M,\theta)\) be a compact strictly pseudoconvex **pseudo‑Einstein** CR manifold with CR dimension \(n\ge2\).  
Let \(Q'_\theta\) denote the CR \(Q'\)-curvature.  For general \(n\), Fefferman–Hirachi’s ambient metric theory classifies scalar pseudohermitian invariants of each weight modulo divergences, and the total \(\int_M Q'_\theta\,dV_\theta\) is known to be a conformal invariant under pseudo‑Einstein changes \(\tilde{\theta}=e^{2u}\theta\).

However, the classification is **fully explicit only in low weights** (and, in dimension \(3\), the weight \(4\) case is particularly delicate):  beyond those, it is unknown whether new, nontrivial, pseudo‑Einstein conformal invariants appear.

---

**Problem (Uniqueness conjecture for CR integral conformal invariants at weight \(2n+2\)):**  
For a strictly pseudoconvex pseudo‑Einstein CR manifold \((M^{2n+1},\theta)\), consider all smooth scalar pseudohermitian invariants \(I(\theta)\) of weight \(2n+2\) constructed as universal complete contractions from the Tanaka–Webster curvature, torsion, and their covariant derivatives.  
Let  
\[
\mathcal{I}[\theta]=\int_M I(\theta)\,dV_\theta.
\]
  
**Question:**  
Assuming that \(\mathcal{I}\) is invariant under pseudo‑Einstein conformal changes (\(\mathcal{I}[\tilde{\theta}]=\mathcal{I}[\theta]\) for every \(\tilde{\theta}=e^{2u}\theta\)), must \(I(\theta)\) coincide (up to a divergence and a constant multiple) with \(Q'_\theta\)?  

Equivalently, is the total CR \(Q'\)-curvature the *unique* nontrivial scalar conformally invariant integral in the pseudo‑Einstein class at weight \(2n+2\)?  

---

Why it is a "good" problem:

1. **Precisely targets an unresolved structural gap.**  
   The Fefferman–Graham–Hirachi framework gives general existence and classification theorems for CR invariants but does not yield a complete uniqueness statement for integral conformal invariants at weight \(2n+2\) under the pseudo‑Einstein restriction.  Explicit confirmation or refutation of uniqueness here remains an open frontier.

2. **Genuinely research‑level.**  
   Resolving uniqueness would likely require a detailed ambient metric or tractor‑calculus analysis in high dimension. It engages the structure of CR invariant theory, parabolic representation theory, and the geometry of the pseudo‑Einstein condition in a manner not settled by existing results.

3. **Conceptually simple, structurally deep.**  
   The statement is short and understandable:  
   > “Is the total \(Q'\)-curvature the *only* global conformal invariant in the pseudo‑Einstein CR class?”  
   Yet it encodes the full algebra of local invariants at a critical weight, akin to the profound uniqueness of the Riemannian \(Q\)-curvature integral.

4. **Connects several core ideas.**  
   It bridges pseudohermitian geometry, integral invariant theory, and conformal analysis, inviting tools from ambient metric expansions, representation theory, and geometric analytic techniques.  Progress would unify these domains.

5. **Bidirectional significance.**  
   A positive answer would show rigidity and a deep parallel with the Riemannian case, confirming \(Q'\) as the fundamental CR curvature scalar.  
   A counterexample would expose previously unknown CR invariants, expanding the current invariant hierarchy and influencing related analytic procedures (such as total \(Q'\)-curvature renormalization and CR anomalies).

This formulation pinpoints a concrete, sharply delimited question—**uniqueness of the \(Q'\)-curvature integral**—that remains structurally motivated, open, and significant for the foundations of pseudohermitian invariant theory.